\chapter{Introdução}

\hl{Falta reescrever a apresentação/preparação inicial deste capitulo aqui
Deve ser algo facil de perceber e de "digerir".}

\section{Enquadramento}
A broncoscopia flexível é um procedimento essencial em Pneumologia, Anestesiologia, Medicina Intensiva e Cirurgia Torácica, cuja execução exige coordenação psicomotora, reconhecimento anatómico, capacidade de navegação endoscópica e tomada de decisão num ambiente clínico variável. Por esse motivo, a simulação tem sido proposta como uma etapa relevante no treino de novos médicos na área, permitindo suavizar a curva de aprendizagem, possibilitar treinos repetidos e aumentar a segurança antes da prática em doentes reais \hl{(será que vale apenas falar já aqui sobre as opções atuais dos vários treinos de simuladores??)} \parencite{Davoudi2008,Kennedy2013,Fielding2014,Nilsson2017,Gopal2018}.

O projeto apresentado neste relatório teve origem numa proposta de trabalho elaborada pela equipa de Pneumologia do Centro Hospitalar Universitário do Algarve \parencite{Barros2023}, que identificou a necessidade de um modelo de simulação de árvore brônquica, impresso em 3D e de baixo custo, para treino de broncoscopia flexível. A partir dessa proposta, o objetivo deste trabalho passou a ser o desenvolvimento de um modelo articulado de uma árvore brônquica, capaz de reproduzir, de forma legítima, os movimentos observáveis durante as broncoscopias, partindo da utilidade já reconhecida dos modelos físicos impressos em 3D e acrescentando-lhes uma componente dinâmica, de modo a aproximar a experiência de treino das condições reais de observação endoscópica \parencite{Pedersen2017,Fu2023}.

O projeto foi dividido em duas secções complementares:
\begin{itemize}
  \item O Modelo Dinâmico, que inclui todo o sistema de controlo e atuação motora, dividido em vários manipuladores delta, bem como o próprio modelo 3D da árvore brônquica;
  \item A HMI, ou interface homem-máquina, alojada num Raspberry Pi 4B com um monitor touchscreen, cuja função é permitir que o utilizador controle o modelo dinâmico de forma simples, intuitiva e natural, através de botões e menus dedicados, sem necessidade de conhecer os comandos da comunicação série entre os dois blocos.
\end{itemize}

\section{Motivação}
Existem vários simuladores comerciais de broncoscopia, contudo estes facilmente podem apresentar custos monetários extremamente elevados, o que limita a sua disponibilidade e acessibilidade em contextos de ensino e treino, razão pela qual a impressão 3D surge como uma forte alternativa, ao permitir criar modelos anatómicos acessíveis, personalizáveis e adaptáveis a diferentes níveis de dificuldade \parencite{Pedersen2017,Fu2023}. Contudo, muitos destes modelos físicos permanecem estáticos, falhando na representação de fenómenos como a deslocação respiratória das vias aéreas, a variação de calibre dos brônquios, as pequenas oscilações induzidas pelo coração e as perturbações rápidas associadas à tosse, resultantes da resposta do corpo humano à inserção de um corpo externo nos brônquios \parencite{Heinbecker1927,Ellis1936,Chen2015,Gunatilaka2020}.

A motivação central deste trabalho consistiu, portanto, na capacidade de combinar a acessibilidade e a versatilidade dos modelos físicos impressos em 3D com a possibilidade de implementação de movimentos anatómicos legítimos, comprovando assim a viabilidade de um simulador dinâmico realista, algo inovador quando comparado aos simuladores estáticos comuns.

\section{Objetivos}
Este projeto pretendia investigar, desenvolver, construir e documentar um protótipo articulado de uma árvore brônquica para treino de videobroncofibroscopia.

Os seus objetivos específicos foram:
\begin{itemize}
  \item estudar requisitos anatómicos, mecânicos e pedagógicos para os simuladores de broncoscopia;
  \item desenvolver um plano e uma estrutura física compatível com impressão 3D e atuação mecânica;
  \item integrar sensores ou atuadores e todo um sistema de eletrónica de controlo flexível, para os vários cenários de treino;
  \item implementar firmware e configurações pré-feitas de controlo para os diferentes perfis de movimento;
  \item desenvolver uma interface, HMI, para um controlo simples e intuitivo do modelo;
  \item definir uma estratégia de testes para avaliação da usabilidade, da repetibilidade e da robustez dos sistemas, bem como verificação do sincronismo dos movimentos em termos de tempo e de espaço percorrido pelos manipuladores robóticos;
  \item documentar as decisões de projeto, limitações e possíveis melhorias futuras.
\end{itemize}

\section{Metodologia}
A metodologia do trabalho combina revisão bibliográfica, definição de requisitos, desenvolvimento iterativo de hardware, desenvolvimento de software embutido e preparação de testes funcionais. \hl{não sei o que escrever mais aqui...}

\section{Organização do relatório}
O relatório está organizado em seis capítulos: 
\begin{itemize}
  \item Capítulo 1: Apresenta o enquadramento, a motivação e os objetivos;
  \item Capítulo 2: Revisão do estado da arte, incluindo treino por simulação, impressão 3D e movimentos observáveis durante a broncoscopia;
  \item Capítulo 3: Descreve o desenvolvimento do hardware;
  \item Capítulo 4: Apresenta estrutura de software, configurações internas, máquinas de estados e todo o sistema por detrás da HMI;
  \item Capítulo 5: Reúne o plano de testes e identifica a informação experimental ainda em falta;
  \item Capítulo 6: Sintetiza as conclusões e melhorias futuras.
\end{itemize}
